\chapter{Algebraic Multigrid}\label{v0ch:20}
\begin{chaptersummary}
Geometric multigrid uses an existing mesh hierarchy; when the main input is a sparse matrix, coarse spaces must be built from algebraic coupling. Classical AMG and smoothed aggregation use strong connections, near-nullspace modes, and PDE information to construct interpolation. Both retain the error directions that smoothing leaves behind.
\end{chaptersummary}
\learninggoals{Explain algebraically smooth error; distinguish classical interpolation from aggregation; record operator complexity and coarse costs.}
\paragraph{Core Question}Given only a sparse matrix, which errors should the coarse level represent?

\section{From Geometric to Algebraic Smoothness}
An error changed little by relaxation is algebraically smooth relative to that smoother. For typical diffusion M-matrices, large negative off-diagonal entries often signal strong coupling. One classical strength criterion is
$-a_{ij}\ge\theta\max_{k\ne i}(-a_{ik})$.
The criterion depends on signs and problem type and does not transfer unchanged to arbitrary nonsymmetric or block systems.
Classical AMG selects coarse variables and builds interpolation through strong connections. Smoothed aggregation first groups variables into tentative interpolation $P_0$, then improves its energy properties using
\[
 P=(I-\omega D^{-1}A)P_0,\qquad A_c=P^TAP
\]
or related operations. Elasticity often requires rigid-body near-null modes; more complex systems may require block information or dedicated strategies.
\conceptbridge{Why does AMG need physical information beyond the nonzero pattern?}{A pattern identifies coupling but not low-energy or poorly relaxed combinations. Coefficient magnitudes, signs, anisotropy, and near-nullspaces determine the directions that interpolation must retain.}{Matrix coupling $\longrightarrow$ difficult error modes $\longrightarrow$ appropriate interpolation space.}
\section{Complexity and Reliable Combinations}
Grid and operator complexities are
\[
 C_g=\frac{\sum_\ell N_\ell}{N_0},\qquad
 C_o=\frac{\sum_\ell\nnz(A_\ell)}{\nnz(A_0)}.
\]
Denser coarse levels may increase setup, memory, and communication even when iteration counts decrease. AMG can be a standalone iteration or an outer preconditioner. Standard CG needs a fixed SPD overall cycle; variable inner strategies generally motivate FGMRES.
\section{An Executable Teaching Model}
The one-dimensional Poisson example combines adjacent aggregation, tentative constant interpolation, one Jacobi smoothing step, and an exact coarse solve into a two-level action. It verifies $P^TAP$, symmetry, and residual correction. Production AMG adds recursive coarsening, robust strength measures, truncation, and parallel coarse-level treatment; performance measurements refer to each complete implementation.\cite{petsc}
\costnote{Report setup, $C_g,C_o$, cycle cost, and outer iterations. Strongly nonsymmetric transport, saddle-point systems, and high-frequency indefinite waves often need specialized preconditioners; diffusion-AMG results do not directly generalize.}
\coderef{sparse_linear_algebra/multilevel.py}
\aiexercise{For seven-point Poisson, construct adjacent-aggregate $P_0$, smoothed $P$, and $P^TAP$. Compare correction of constant and oscillatory errors.}{Vary coefficient jumps and aggregation. Record complexity, setup, and PCG true residuals, checking symmetry and positive definiteness first.}
\studyroute{Ruge--St\"uben interpolation formulas and smoothed-aggregation convergence theory are omitted. See \cite{saad,briggs,petsc}.}{algebraically smooth error; strength of connection; classical AMG; smoothed aggregation; near nullspace}
